\chapter{Heat and Work}

\begin{definition}[Heat]
Heat is any spontaneous flow of energy from one object to another caused by a difference in temperature.
\end{definition}

\begin{center}[DIAGRAM: Two blocks in thermal contact --- one hot, one cold --- with wavy arrows indicating heat flowing from the hotter to the cooler block]\end{center}

\begin{definition}[Work]
Work is any transfer of energy into or out of a system other than heat. Common examples include pushing a piston, stirring a fluid, or driving a current through a resistor.
\end{definition}

\section{First Law of Thermodynamics}

\begin{theorem}[First Law]
\[
    \Delta U = Q + W,
\]
where $U$ is the total internal energy of the system, $Q$ is the heat added to the system, and $W$ is the work done on the system.
\end{theorem}

Heat can be transferred between systems by three mechanisms. \textbf{Conduction} is the transfer of thermal energy through direct molecular contact. \textbf{Convection} is driven by the bulk motion of a fluid: warmer material expands, becomes less dense, and rises, continuously transporting energy. \textbf{Radiation} is the emission and absorption of electromagnetic waves, which requires no material medium.

\subsection{Units of Heat}

The \textbf{calorie} is defined as the amount of heat required to raise the temperature of one gram of water by one degree Celsius. The mechanical equivalent of heat,
\[
    1\,\text{cal} \approx 4.2\,\text{J},
\]
establishes that heat and work are both forms of energy transfer, measured in the same units.

\section{Compression Work}

\begin{center}[DIAGRAM: A piston of area $A$ being pushed inward by a force $F$ through a displacement $\Delta x$, compressing gas in a cylinder. The volume change is $\Delta V = -\Delta x \cdot A$.]\end{center}

The work done on a gas by a moving piston follows from $W = \int \vec{F} \cdot d\vec{x}$. In the limit of \textbf{quasistatic compression} --- compression slow enough that the pressure remains spatially uniform throughout the gas at every instant --- the force exerted on the piston equals $F = PA$. Since moving the piston inward by $dx$ decreases the volume by $dV = -A\,dx$, the work done on the gas is

\begin{formula}[Compression Work]
\[
    W = -\int_{V_i}^{V_f} P(V)\,dV.
\]
\end{formula}

When the gas is compressed ($V_f < V_i$) the integral is negative, so $W > 0$: positive work is done on the gas, as expected.

\subsection{Isothermal Compression}

Isothermal compression is carried out so slowly that the gas remains in thermal contact with its surroundings and its temperature stays constant throughout. With $T$ fixed, the ideal gas law gives $P = NkT/V$, so

\begin{formula}[Isothermal Work]
\[
    W = -\int_{V_i}^{V_f} \frac{NkT}{V}\,dV = -NkT\ln\frac{V_f}{V_i} = NkT\ln\frac{V_i}{V_f}.
\]
\end{formula}

Since the temperature is constant, the internal energy of an ideal gas is unchanged: $\Delta U = \frac{f}{2}Nk\Delta T = 0$. The first law then gives $Q = -W$, so all the work done on the gas flows out as heat:

\[
    Q_{\text{isothermal}} = -NkT\ln\frac{V_i}{V_f} = NkT\ln\frac{V_f}{V_i}.
\]

\subsection{Adiabatic Compression}

Adiabatic compression is carried out so rapidly that no heat has time to enter or leave the system: $Q = 0$. The compression is still assumed to be quasistatic, so the work formula above still applies. The first law reduces to $dU = dW$, giving

\[
    \frac{f}{2}Nk\,dT = -P\,dV.
\]

Dividing both sides by $NkT = PV$ separates variables:

\[
    \frac{f}{2}\frac{dT}{T} = -\frac{dV}{V}.
\]

Integrating from initial to final state:

\[
    \frac{f}{2}\ln\frac{T_f}{T_i} = -\ln\frac{V_f}{V_i},
\]

which exponentiates to the adiabatic constraint:

\begin{formula}[Adiabatic Constraint (T, V)]
\[
    V_f T_f^{f/2} = V_i T_i^{f/2}.
\]
\end{formula}

To express the constraint in terms of pressure and volume, substitute $T = PV/(Nk)$:

\[
    V_f\!\left(\frac{P_f V_f}{Nk}\right)^{f/2} = V_i\!\left(\frac{P_i V_i}{Nk}\right)^{f/2}
    \implies V_f^{f/2+1}\,P_f^{f/2} = V_i^{f/2+1}\,P_i^{f/2}.
\]

Defining the \textbf{adiabatic index} $\gamma \equiv (f+2)/f$, so that $f/2 + 1 = f\gamma/2$, this simplifies to:

\begin{formula}[Adiabatic Constraint (P, V)]
\[
    V_f^\gamma P_f = V_i^\gamma P_i, \qquad \gamma = \frac{f+2}{f}.
\]
\end{formula}

The quantity $PV^\gamma$ is conserved along an adiabat. Since $\gamma > 1$, the adiabat $P \propto V^{-\gamma}$ falls more steeply than the isotherm $P \propto V^{-1}$ in a $P$--$V$ diagram: when you compress a gas adiabatically, the pressure rises faster because the temperature also rises.

\begin{center}[DIAGRAM: $P$--$V$ diagram showing an isotherm (shallower curve, $P \propto 1/V$) and an adiabat (steeper curve) both passing through the same initial point $(V_i, P_i)$ and ending at the same final volume $V_f$, with the adiabat reaching a higher final pressure.]\end{center}

\section*{Problem 1.38}

A bubble of gas expands as it rises through water. Compare the final volumes under isothermal and adiabatic expansion for the same pressure drop.

\textit{Isothermal:} $PV = \text{const}$, so $P_f V_f = P_i V_i$, giving $P_f/P_i = V_i/V_f$.

\textit{Adiabatic:} $PV^\gamma = \text{const}$, so
\[
    \frac{P_f}{P_i} = \left(\frac{V_i}{V_f}\right)^\gamma
    \implies \left(\frac{P_f}{P_i}\right)^{1/\gamma} = \frac{V_i}{V_f}.
\]

For the same pressure ratio $P_f/P_i < 1$, we compare the two expressions for $V_i/V_f$. Since $\gamma > 1$ implies $1/\gamma < 1$, raising a number less than one to a smaller power brings it closer to one, so

\[
    \left(\frac{P_f}{P_i}\right)^{1/\gamma} < \frac{P_f}{P_i}
    \implies \left(\frac{V_i}{V_f}\right)_\text{adiabatic} < \left(\frac{V_i}{V_f}\right)_\text{isothermal}
    \implies \boxed{V_{f,\,\text{adiabatic}} > V_{f,\,\text{isothermal}}}.
\]

This makes physical sense: during adiabatic expansion the gas cools, which lowers the pressure for a given volume, causing it to expand further to reach the same final pressure. Equivalently, the isothermal bubble has had time to absorb heat from the surroundings and thus expands less.

\chapter{Heat Capacity and Latent Heat}

\section{Heat Capacities}

\begin{definition}[Heat Capacity]
The \textbf{heat capacity} $C$ of an object is the amount of heat required to raise its temperature by $1\,\text{K}$:
\[
    C = \frac{Q}{\Delta T}.
\]
The \textbf{specific heat capacity} $c = C/m$ normalises by mass.
\end{definition}

The heat capacity is not uniquely defined without specifying the conditions under which heat is added, since $Q$ depends on whether the volume or pressure is held fixed. These two cases give different, well-defined quantities.

\subsection{Heat Capacity at Constant Volume}

When volume is held fixed, no expansion work is done ($W = -P\Delta V = 0$), so the first law gives $\Delta U = Q$. The heat capacity at constant volume is therefore

\begin{formula}[Heat Capacity at Constant Volume]
\[
    C_V = \left(\frac{\partial U}{\partial T}\right)_{\!V}.
\]
\end{formula}

For a system whose energy is stored entirely in quadratic degrees of freedom, the equipartition theorem gives $U = \frac{f}{2}NkT$, and differentiating yields

\begin{theorem}[Dulong--Petit Rule]
\[
    C_V = \frac{f}{2}Nk = \frac{f}{2}nR.
\]
The heat capacity is proportional to the number of active degrees of freedom per molecule.
\end{theorem}

This result is remarkably simple: each degree of freedom contributes $\frac{1}{2}Nk$ to the heat capacity. In practice, $f$ is not constant with temperature --- various degrees of freedom (vibrational modes in particular) ``freeze out'' at low temperatures as the spacing between quantum energy levels becomes large compared to $kT$, reducing the effective $f$ and hence $C_V$.

\subsection{Heat Capacity at Constant Pressure}

At constant pressure, the gas is free to expand when heated, so it does work on its surroundings. The heat capacity at constant pressure is therefore larger than $C_V$:

\[
    C_P = \left(\frac{Q}{\Delta T}\right)_{\!P}
    = \frac{\Delta U - W}{\Delta T}
    = \frac{\Delta U + P\Delta V}{\Delta T}
    = \left(\frac{\partial U}{\partial T}\right)_{\!P} + P\left(\frac{\partial V}{\partial T}\right)_{\!P}.
\]

For an ideal gas, $V = NkT/P$, so $(\partial V/\partial T)_P = Nk/P$, giving $P(\partial V/\partial T)_P = Nk$. Since $(\partial U/\partial T)_P = C_V$ for an ideal gas (internal energy depends only on temperature), we obtain

\begin{formula}[$C_P$ for an Ideal Gas]
\[
    C_P = C_V + Nk = C_V + nR.
\]
\end{formula}

The heat capacity at constant pressure exceeds that at constant volume by exactly $nR$, regardless of the molecular structure of the gas. The extra energy goes into the work of expansion rather than into raising the temperature.

\section{Latent Heat}

During a phase transition (melting, boiling), a system absorbs heat without any change in temperature. The ordinary definition $C = Q/\Delta T$ would give an infinite heat capacity, which signals that temperature is simply the wrong variable here.

\begin{definition}[Latent Heat]
The \textbf{latent heat} $L$ is the total heat required to completely convert a substance from one phase to another at constant pressure. The \textbf{specific latent heat} is $\ell = L/m$, measured in $\text{J/kg}$ or $\text{J/g}$.
\end{definition}

The energy absorbed during a phase transition goes into breaking intermolecular bonds rather than increasing the kinetic energy of molecules, which is why the temperature remains constant throughout. Typical values:
\begin{itemize}
    \item Melting ice: $\ell = 333\,\text{J/g}$ ($80\,\text{cal/g}$)
    \item Boiling water: $\ell = 2260\,\text{J/g}$ ($540\,\text{cal/g}$)
\end{itemize}

\section*{Problems}

\subsection*{Problem 1.41}

A metal block is dropped into a cup of water. Since the system is thermally isolated, $\Delta U_\text{total} = 0$, which means the heat lost by the metal equals the heat gained by the water:
\[
    m_m C_m \,|\Delta T_m| = m_w C_w \,|\Delta T_w|.
\]

\begin{enumerate}[label=\alph*)]
    \item The heat gained by the water: $Q_w = m_w C_w \Delta T_w = 250 \times 4 \times 4 = 4000\,\text{J}$.
    \item By conservation of energy, the metal lost the same amount: $4000\,\text{J}$.
    \item The heat capacity of the metal block: $m_m C_m = \dfrac{4000\,\text{J}}{78\,\text{K}} \approx 51.3\,\text{J/K}$.
    \item The specific heat capacity: $C_m = \dfrac{51.3}{100\,\text{g}} \approx 0.51\,\text{J/(g\,K)}$.
\end{enumerate}

\subsection*{Problem 1.43}

The measured molar heat capacity of water is $C_V/\text{molecule} = 18 \times 4.186\,\text{J} / (6.022 \times 10^{23}) \approx 9.07\,k$. From the equipartition theorem, $C_V/\text{molecule} = fk/2$, so

\[
    f \approx 2 \times 9.07 \approx 18.
\]

This large value (compared to $f=3$ for a monatomic gas) reflects the many vibrational and rotational modes active in a liquid water molecule at room temperature.

\medskip

The problem also introduces two thermodynamic response functions. The \textbf{thermal expansion coefficient} and \textbf{isothermal compressibility} are defined by
\[
    \beta = \frac{1}{V}\left(\frac{\partial V}{\partial T}\right)_{\!P}, \qquad
    \kappa_T = -\frac{1}{V}\left(\frac{\partial V}{\partial P}\right)_{\!T}.
\]

\begin{enumerate}[label=\alph*)]
    \item A temperature increase $dT$ at constant pressure produces a volume change $dV = \beta V\,dT$.

    \item A pressure increase $dP$ at constant temperature produces $dV = -\kappa_T V\,dP$.

    \item Along a path of constant volume, the two effects must cancel: $\beta V\,dT = \kappa_T V\,dP$, giving the useful identity
    \[
        \left(\frac{\partial P}{\partial T}\right)_{\!V} = \frac{\beta}{\kappa_T}
        = -\frac{(\partial V/\partial T)_P}{(\partial V/\partial P)_T}.
    \]

    \item For an ideal gas, $V = nRT/P$, so
    \[
        \beta = \frac{1}{V}\cdot\frac{nR}{P} = \frac{1}{T}, \qquad
        \kappa_T = \frac{1}{V}\cdot\frac{nRT}{P^2} = \frac{1}{P}.
    \]
    Therefore $(\partial P/\partial T)_V = \beta/\kappa_T = P/T$, which agrees with differentiating $PV = nRT$ at constant $V$. \checkmark

    \item A real-gas pressure rise along a constant-volume path is $\partial P = (\beta/\kappa_T)\,\partial T$.
\end{enumerate}

\subsection*{Problem 1.48}

Estimate how long a glacier of cross-section $2\,\text{m}\times 1\,\text{m}$ and height $100\,\text{m}$, composed of $50\%$ ice and $50\%$ air by volume, would take to melt.

\begin{center}[DIAGRAM: Rectangular block $2\,\text{m} \times 1\,\text{m}$ base and $100\,\text{m}$ height.]\end{center}

The volume of the block is $2\,\text{m}^3$ per metre of height, giving a total of $200\,\text{m}^3$, of which half is ice: $100\,\text{m}^3$ of ice with mass $917\,\text{kg/m}^3 \times 100\,\text{m}^3 = 9.17 \times 10^4\,\text{kg} = 9.17 \times 10^7\,\text{g}$.

The energy required to melt all the ice:
\[
    Q = \ell \cdot m = 333\,\text{J/g} \times 9.17\times10^7\,\text{g} \approx 3.1\times10^{10}\,\text{J}.
\]

Estimating the absorbed solar power as roughly $1\,\text{kW/m}^2$ over the $2\,\text{m}^2$ top surface gives $\sim 2\,\text{kW}$, so
\[
    t \approx \frac{3.1\times10^{10}\,\text{J}}{2\times10^3\,\text{W}} \approx 1.5\times10^7\,\text{s}.
\]

Converting: $1.5\times10^7\,\text{s} \div (6\times10^5\,\text{s/week}) \approx 4.5\,\text{weeks} \approx 5\,\text{weeks}$.
